% !PRINTED 332
\weberpage{332}
Aus dem Ausdruck (1) können wir eine \gesperrt{canonische Darstellung}
 für die cubische Form $f$ herleiten:
Wir bezeichnen mit $\varepsilon$ eine imaginäre dritte Einheitswurzel
und führen zwei Variable $z_1, z_2$ ein, die durch die Gleichungen

\weq{2}{\begin{aligned} a_1\,x_1 &= \varepsilon\,z_1 + \varepsilon^2\,z_2 - \tfrac{1}{3}\,a_3\,x_3 \\ a_2\,x_2 &= \varepsilon^2\,z_1 + \varepsilon\,z_2 - \tfrac{1}{3}\,a_3\,x_3 \end{aligned}}

definirt sind, woraus noch folgt:

\weq{3}{-\,(a_1\,x_1 + a_2\,x_2 + a_3\,x_3) = z_1 + z_2 - \tfrac{1}{3}\,a_3\,x_3.}

Die Variablen $z_1, z_2$ sind durch (2) als lineare Functionen
von $x_1, x_2, x_3$ bestimmt, und die linearen Functionen $z_1, z_2, x_3$
sind als Functionen von $x_1, x_2, x_3$ linear unabhängig. Wenn wir
(2) und (3) in (1) einführen, so ergiebt sich

\[
-\,a_1\,a_2\,f = z_1^3 + z_2^3 - \left(a_1\,a_2\,h + \frac{a_3^3}{27}\right)x_3^3 + a_3\,z_1\,z_2\,x_3.
\]

Diesem Ausdrucke giebt man eine mehr symmetrische Gestalt
, indem man für $z_1, z_2, x_3$ drei neue lineare Functionen
$y_1, y_2, y_3$ einführt, die sich von diesen nur um constante Factoren
unterscheiden:

\[
f = \varphi\,(y_1,y_2,y_3) = \alpha_1\,y_1^3 + \alpha_2\,y_2^3 + \alpha_3\,y_3^3 + 6\,m\,y_1\,y_2\,y_3,
\]

worin $\alpha_1, \alpha_2, \alpha_3, m$ Constanten bedeuten.

Wenn kein singulärer Punkt vorhanden ist, so können die
Coëfficienten $\alpha_1, \alpha_2, \alpha_3$ nicht verschwinden, und man kann sie
ohne Beschränkung der Allgemeinheit einander gleich, etwa $=h$,
annehmen. Man könnte sie sogar $=1$ setzen, was aber, um die
Homogenëität aufrecht zu erhalten, zunächst besser nicht geschieht.
Wir haben also die \gesperrt{canonische Form} für die ternäre
cubische Form:

\weq{4}{\varphi\,(y_1,y_2,y_3) = h\,(y_1^3+y_2^3+y_3^3) + 6\,m\,y_1\,y_2\,y_3.}

Nach unserem Ausgangspunkte war die Linie $x_3$, oder was
dasselbe ist, $y_3$ die Verbindungslinie dreier Inflexionspunkte.
Solcher Linien giebt es aber zwölf, und wir können daher die
Linie $y_3$ auf zwölf Arten wählen. Setzen wir aber $y_1$ oder $y_2$ an
Stelle von $y_3$, so bekommen wir dieselbe Darstellung in der
canonischen Form, und es giebt also nur vier wesentlich verschiedene
 Arten dieser Darstellung (abgesehen von dem willkürlichen
 $h$), d.\,h. vier Arten, die Linien $y_1, y_2, y_3$ zu bestimmen. Wir
haben daher den Satz:
%CONT>>
